\section{Optuna}

\textbf{Optuna} es un framework de optimización de hiperparámetros de código abierto que permite a los desarrolladores y científicos de datos encontrar automáticamente los mejores hiperparámetros para sus modelos de aprendizaje automático. Optuna utiliza técnicas de optimización bayesiana para explorar el espacio de hiperparámetros de manera eficiente, lo que puede mejorar significativamente el rendimiento de los modelos.

A diferencia de Grid Search, que realiza una búsqueda exhaustiva y costosa, o Random Search, que no aprende de iteraciones pasadas, Optuna utiliza un enfoque bayesiano basado en el Tree-structured Parzen Estimator (TPE).

El algoritmo TPE construye un modelo probabilístico de la función objetivo. En lugar de probar combinaciones al azar, predice qué regiones del espacio de búsqueda tienen más probabilidades de mejorar los resultados actuales.

Otra de las ventajas de este framework es que nos permite establecer una \textbf{optimización multiobjetivo}, lo que es especialmente útil cuando se desea optimizar varias métricas al mismo tiempo, como la precisión y eficiencia del modelo.

Optuna identifica un conjunto de soluciones que se encuentran en la frontera de Pareto, lo que significa que no hay otra solución que sea mejor en todas las métricas. Esto permite a los usuarios elegir entre diferentes compromisos según sus necesidades específicas, por lo que no presentamos una única solución.

\begin{figure}[H]
    \centering
    \includegraphics[width=0.8\textwidth]{images/OPTUNA.png}
    \caption{Visualización de la optimización con Optuna}
    \label{fig:optuna}
\end{figure}

El procedimiento general aplicado a cada modelo se resume en el Algoritmo~\ref{alg:optuna_cv_multiobjetivo}. Aunque cada arquitectura tiene un espacio de hiperparámetros específico, la estrategia de evaluación se mantiene constante: cada configuración se evalúa mediante validación cruzada estratificada y se optimiza simultáneamente el rendimiento predictivo y el coste de inferencia.

\begin{algorithm}[H]
\caption{Optimización multiobjetivo con validación cruzada estratificada}
\label{alg:optuna_cv_multiobjetivo}
\begin{algorithmic}[1]
\Require Conjunto de entrenamiento completo $D_{train} = (X, y)$, modelo base $M$, espacio de hiperparámetros $\mathcal{H}$, número de ensayos $N$
\Ensure Conjunto de configuraciones evaluadas y frontera de Pareto

\State Crear un estudio de Optuna con dos objetivos:
\Statex \hspace{\algorithmicindent} \textbf{maximizar $F1$ y minimizar la latencia media de inferencia}
\State Definir una validación cruzada estratificada con $K=3$ particiones

\For{cada ensayo $t = 1, \dots, N$}
    \State Seleccionar una configuración de hiperparámetros $h_t \in \mathcal{H}$
    \State Inicializar las listas $F1\_scores$ y $Latencias$

    \For{cada partición $(D_{train}^{(k)}, D_{val}^{(k)})$ de la validación cruzada}
        \State Instanciar el modelo $M_t$ con los hiperparámetros $h_t$
        \State Entrenar $M_t$ sobre $D_{train}^{(k)}$
        \State Obtener las predicciones sobre $D_{val}^{(k)}$
        \State Calcular el $F1$ binario de la clase maliciosa
        \State Añadir el valor obtenido a $F1\_scores$

        \State Seleccionar un subconjunto de validación para medir latencia
        \State Realizar una predicción inicial de calentamiento
        \State Medir varias repeticiones del tiempo de inferencia
        \State Calcular la latencia media por muestra en milisegundos
        \State Añadir la latencia obtenida a $Latencias$
    \EndFor

    \State Calcular $F1_{medio}$ como la media de $F1\_scores$
    \State Calcular $Latencia_{media}$ como la media de $Latencias$
    \State Devolver a Optuna el par $(F1_{medio}, Latencia_{media})$
\EndFor

\State Obtener la frontera de Pareto a partir de todos los ensayos evaluados

\end{algorithmic}
\end{algorithm}
